\begin{proposition}\label{prop:nonic57-quotient9-partial}
Put $X=X_1(57)$ and $H=\langle4\rangle\subset
(\mathbb Z/57\mathbb Z)^*/\{\pm1\}$.
The following are necessary conditions for a noncuspidal closed
point of degree $9$ on $X$.
At $2$, its noncuspidal reduction has degree $6$ or $8$.
On the genus-$9$ quotient by $H$, there is one remaining normed
image in degree $6$ and two remaining images in degree $8$.
If its reduction is entirely cuspidal, it is the reduction of a
rational cuspidal divisor.
\end{proposition}
\begin{proof}
We use the rational $q$-expansion model $Y$ of the quotient.
This is the dual moduli convention: a point level structure
is sent to the quotient elliptic curve with its dual multiplicative
level structure. Fricke interchanges the displayed $j$-widths
$w$ and $57/w$. The map from $X$ to this model is defined over
$\mathbb Q$ and has degree $9$.

The genus of $Y$ is $9$. The exact modular-symbol decomposition
has dimensions $1,1,1,1,4$, with the level-$19$ elliptic factor
occurring twice. Only the first level-$57$ elliptic factor has
zero central $L$-value. The involution $\sigma=\langle5\rangle$
acts as $-1$ precisely on the four-dimensional factor, whose
exact central ratio is $32/49$.
Consequently $(\sigma-1)J_Y(\mathbb Q)$ is finite.
The associated abelian fourfold has $49$ points over $\mathbb F_2$
and $441$ points over $\mathbb F_5$. Its rational torsion order
therefore divides $49$.

We construct a saturated canonical model, check all $21$
quadrics on its normalized genus-$9$ plane model, and recover
the $j$-map from both the $E_4^3$ and $\Delta$ section conditions.
The map has degree $160$. Its finite cusp signature is
$1^4 2^2$, whereas its global cusp signature is $1^2 2^3$.
The proven finite Jacobian is
\[
 J_Y(\mathbb F_2)=[105,210].
\]
All finite cusp classes generate $[21,210]$.
Their images under $\alpha=\sigma-1$ generate $[7,7]$;
this is the entire finite image of $\alpha$, since the image
abelian fourfold has order $49$ over $\mathbb F_2$.

In the dual convention the Eichler--Shimura annihilator at $5$ is
\[
 A_5=T_5-\langle5\rangle-5.
\]
It kills the rational finite image of $\alpha$.
Indeed that image has rank zero and order prime to $5$.
We compute its action at the cusps by the dual boundary formula
\[
 T_5=\langle5\rangle\mathrm{Frob}_5+
 5\mathrm{Frob}_5^{-1}.
\]
In the exact finite cusp order, the widths, diamond permutation
and Galois permutation are respectively
\[
 (1,19,57,3,57,1),\quad
 (6,2,5,4,3,1),\quad(1,2,5,4,3,6).
\]
The induced operator respects every relation of the finite
$\alpha$-image, and its kernel is a subgroup $K\simeq C_7$.
Thus the reduction of the rational $\alpha$-image is contained
in $K$.

For $d\leq5$, the Hasse bound gives
$\#E(\mathbb F_{2^d})<57$, so no noncuspidal source place
has degree at most $5$. An independent enumeration of every
finite-field $j$-invariant and every twist gives the qualified
$j$-polynomial counts
\[
\begin{array}{c|rrrr}
 d&6&7&8&9\\ \hline
 \text{qualified twists}&2&14&8&1\\
 \text{closed }j\text{-polynomials}&1&2&1&1.
\end{array}
\]
Both entries in degrees $6$ and $9$ have $j=0$.
Since $57$ is squarefree, divisibility of the curve order by
$57$ is equivalent to existence of a point of order $57$.
It is preserved by isogeny, so this same census contains the
$j$-values on the dual quotient model.

The source cusp signatures globally and at $2$ are
\[
 1^{18}2^9 18^2,\qquad 1^{18}2^9 9^2 18.
\]
All cusp places of degree at most $3$ are globally intact.
They are the point-model width-$57$ rational cusps and
width-$19$ quadratic cusps. Their images have dual widths
$1$ and $3$ and are rational closed cusp divisors.
A noncuspidal support of degree $d\geq6$ therefore has a
rational cuspidal complement of degree at most $3$.
After normalization at a rational cusp, its normed $\alpha$-class
must belong to $K$.

We enumerate every closed place of $Y$ of degree at most $9$.
Their counts are
\[
 (6,9,0,2,0,0,24,24,64).
\]
Independent canonical coordinates account for every place and
agree with the exact Frobenius polynomial.
For every $d=6,7,8,9$, we retain every place of degree $e\mid d$
whose $j$-polynomial occurs in the qualified degree-$d$ census,
and apply the norm multiplicity $d/e$.
This containing set has $23$ cases. Testing their normalized
$\alpha$-classes against $K$ leaves exactly
\[
 (d,e,d/e)=(6,2,3),(8,8,1),(8,8,1).
\]
The first image has $j=0$ and is fixed by $\sigma$.
The other two are interchanged by $\sigma$ and lie above
\[
 x^8+x^7+x^6+x^3+x^2+x+1.
\]
Thus all noncuspidal degree-$7$ and degree-$9$ reduction
possibilities are excluded.

Finally the only non-global source cusp places of degree at
most $9$ are the two nonic places splitting one global
cusp orbit of degree $18$. Each maps, with norm $9$, to one of
the two dual width-$57$ cusps.
For both choices, the normalized $\alpha$-class lies outside $K$.
These two entirely cuspidal reductions are also excluded.
Every other entirely cuspidal reduction is globally intact,
as claimed.
\end{proof}

The three surviving images and the globally intact cuspidal
reductions remain to be excluded or lifted. This proposition
does not determine whether $C_{57}$ occurs in degree $9$.
